\section{The refutation register}
\label{sec:register}

A refuted claim is evidence that it was tried. This section reproduces
the register of every claim this program made, or inherited from its own
sources, and then lost --- together with the exact thing that killed it.

The register carries \ContradictionTotal{} entries in three states.
\ContraFalsified{} are \emph{falsified}: a counterexample or an exact
computation contradicted them. \ContraSuperseded{} are
\emph{superseded}: replaced by a sharper statement. \ContraResolved{} are
\emph{resolved}: the apparent conflict dissolved, in most cases because
the two sides turned out to be differently anchored quantities rather
than competing values for the same quantity.

Three features of how this register is maintained matter to a reader
deciding how much to trust the rest of the document.

First, a dispute closes by derivation or it does not close. Both values
stay recorded side by side, and no code path is permitted to pick one,
average them, or prefer whichever is an exact rational because exactness
reads as authority. The fourth-order off-axis coefficient closed that way:
every cluster of the rotation record reproduces the historical pipeline's
own number, and the single cluster that does not is explained ordering by
ordering (\S\ref{sec:fourth}). Its two recorded values remain in the
registry with their verdicts attached.

Second, a resolved entry is not a win. The largest class here is
resolutions in which a contradiction turned out to be a naming
collision --- two differently anchored coordinates both called $m_4$, for
instance. Those cost real time and produced no mathematics, and the
register's response was to forbid the ambiguous name rather than to
record a victory.

Third, this program has retracted a mechanism of its own. The
\emph{tier collapse} --- the observation that two of the four
fourth-order shape coefficients vanish at every solved rank --- was
proposed here as a kinematic degree bound with a falsifiable
prediction, and the mechanism was withdrawn a fortnight later when one
projection turned out to be uncounted (\S\ref{sec:fourth}). The
withdrawal is a numbered decision record rather than a silent edit,
because the next attempt at that mechanism needs to know which step
failed. The reformulation survives; the mechanism and its prediction do
not.

That is a different object from the degree bound of
\S\ref{sec:symbolic}, which was never retracted --- it was made
unnecessary. The two are easy to confuse and this document keeps them
apart deliberately.

\medskip
% Generated by paper/master/generate_apparatus.py.  Do not edit.
\subsection*{Falsified (4)}
\addcontentsline{toc}{subsection}{Falsified}
\noindent Falsified --- killed by a counterexample or an exact computation. These are the four claims this program asserted, or inherited, and then lost.\par\smallskip
\begin{itemize}[leftmargin=1.4em]
  \item \texttt{C6} \textbf{F-2 mobility rule.} The promotion r\_\allowbreak{}physical = w\_\allowbreak{}min - 2 is falsified by the pentagonal isotropic cap hop at r=4 with w\_\allowbreak{}min-2=5. The scoped bound r >= w\_\allowbreak{}min - 2 survives, with survival gates.
  \item \texttt{C7} \textbf{Stranded-flux zero backend.} Claimed Haar zero for all 20 pentagonal histories refuted by exact Weingarten/network contractions (8/8).
  \item \texttt{C18} \textbf{Center-only circuits are dynamically dark.} Refuted by the exact fifth-order determinant correction delta c\_\allowbreak{}5,det != 0
  \item \texttt{C19} \textbf{Global fixed-window firewall.} Disproved by the Bernoulli no-go; replaced by the rooted conditional calculus
\end{itemize}

\subsection*{Superseded (3)}
\addcontentsline{toc}{subsection}{Superseded}
\noindent Superseded --- replaced by a sharper statement. The original is not wrong so much as no longer the best available.\par\smallskip
\begin{itemize}[leftmargin=1.4em]
  \item \texttt{C5} \textbf{String-tension sign at O(u\textasciicircum{}3).} -61/408 (the earlier +61/408 was an odd-order convention error)
  \item \texttt{C11} \textbf{Compact-gap constants.} Corrections: c\_\allowbreak{}0 = -3(N\textasciicircum{}2-3)/(16N); radial moment difference 3N\textasciicircum{}2+1; dilation eps\textasciicircum{}4 = N/(2*beta);
  \item \texttt{C12} \textbf{Curvature misreads.} Two-stage. The kappa values are RADIAL DIRECTIONAL curvatures, not Hessian entries at Gamma; a genuine Hessian needs beta = 2*alpha, violated by the historical kernel.
\end{itemize}

\subsection*{Resolved (15)}
\addcontentsline{toc}{subsection}{Resolved}
\noindent Resolved --- the collision was dissolved. In most of these the two sides turned out to be differently anchored quantities rather than competing values.\par\smallskip
\begin{itemize}[leftmargin=1.4em]
  \item \texttt{C1} \textbf{Fourth-order rest scalar --- dissolved, an anchoring distinction.} NOT a contradiction. q\_\allowbreak{}band\textasciicircum{}(4) and m\_\allowbreak{}Gamma\textasciicircum{}(4) are not two competing estimates of the same scalar coordinate: the first is a band-kernel anchor, the second a vacuum-subtracted physical Gamma-point coefficient. They are differently anchored coordinates related by a translation-local scalar shift, and such a shift cannot change the centered operator, its eigenvectors, the SOS factorization, the mobility coefficients, or the bandwidth.
  \item \texttt{C2} \textbf{Fourth-order off-axis coefficient C\_\allowbreak{}shp --- resolved by derivation 2026-09-04.} C\_\allowbreak{}shp = C\_\allowbreak{}SHP\_\allowbreak{}HISTORICAL + 25/1024 = -13035490122347/550663802582400 = -0.0236723, the cluster-assembled coefficient read in the kernel's own basis (ADR 0024). Both recorded kernels are wrong on the rotation amplitude: the historical by the sixteen orderings of the adjacent-face cube completion its own Stage-3I ledger lacks (8 of 24, each with the right weight; the shortfall is exactly 25/512 on rho), the cold dump on u, pi and rho alike.
  \item \texttt{C3} \textbf{Status of the fourth-order "theorem".} Newest status controls. The \S{}5.2 values are sealed; everything else is C1/C2.
  \item \texttt{C4} \textbf{Coupling conventions.} The archived Y = 2*beta/3 = 4u line is a definition/label error, not a rescaling rule. Printed coefficients were already in u = beta\_\allowbreak{}lat/6. Never apply 4**n factors to them.
  \item \texttt{C8} \textbf{Pentagonal degenerate space.} E\_\allowbreak{}cap = 10/3 != E\_\allowbreak{}side = 8/3; choose the physical degenerate eigenspace first
  \item \texttt{C9} \textbf{SU(4) exceptional correction.} Only the flat-branch eigenvalue identity holds, not scalarity of the matrix
  \item \texttt{C10} \textbf{Determinant sectors shift only q\_\allowbreak{}N.} True for N>=4, false at SU(3): Delta\_\allowbreak{}beta\_\allowbreak{}3 = -25/64 and Delta\_\allowbreak{}q\_\allowbreak{}3 = -16863189551/76406976000.
  \item \texttt{C13} \textbf{C-even second-order hopping.} -481/612 superseded by -11/306 (the vacuum-mediated route was omitted)
  \item \texttt{C14} \textbf{Betti "+2".} Reinterpreted as b\_\allowbreak{}2 - b\_\allowbreak{}3; the open box has no extra states
  \item \texttt{C15} \textbf{Tetrahedral local coefficient.} Re-derived 2026-08-21 (tetrahedral Haar-resolvent suite): the asserted -8/(N(N\textasciicircum{}2-1)) is exact at symbolic N. The derivation reconstructs the primitive law's three stated ingredients --- 1/N per merge from the fundamental-pair Haar moment (independent of shared-path length, derived not assumed), the certified electric convention E(L) = L*C\_\allowbreak{}F/2, and exhaustive temporal orderings --- and reproduces every certified or printed instance first (sealed-core cube -5/48 with the (-8,-8,-4) temporal classes, prism 64 with its cap-sector 24 dissolved as a sector\,\ldots
  \item \texttt{C16} \textbf{SU(5)/SU(6) fourth-order completeness.} The payloads were shipped after all, and are now machine-checked (restored-payloads suite, 2026-08-21): the SU(6) determinant certificate re-derives Delta q\_\allowbreak{}6 = 6/343 exactly, momentum-independent with A, B, and the bandwidth unchanged, and the SU(5) stage-1 summary pins the 895,524-pair scan with zero determinant sectors. Two upstream word manifests remain absent (their SHAs are recorded in the certificates), so the evidence level is output-certified, not cold-reproduced.
  \item \texttt{C17} \textbf{Fixed-rank holdout overclaim.} Corrected: q\_\allowbreak{}N matches stored values for N=7..18, but full kernels only at N=7,8; the A/B sample ledger remains open.
  \item \texttt{C20} \textbf{RUN15 internal display artifact.} Cosmetic; the gate value is the intended exact one
  \item \texttt{C21} \textbf{Monte Carlo gate integrity.} next14.json's "23/23 gates" is downgraded: one gate is a literal truth value, the JSON is not source-hash bound (non-RFC NaN tokens), and no raw ensemble exists. Numbers stand as structured numerical evidence; the certificate claim does not.
  \item \texttt{C22} \textbf{Gate-85 equality.} Not a numerical contradiction but a status correction: the final diagonal shift was chosen to produce exactly that equality, so gate 85 certifies internal bookkeeping, not independent agreement.
\end{itemize}


% status histogram: {'proven': 140, 'disputed': 1, 'falsified': 1, 'conditional': 1}
% evidence histogram: {'analytic': 141, 'record-backed': 1, 'numerical': 1}


\subsection{Claims withdrawn by this program}

Distinct from the register above, the following are claims this
repository itself advanced and then withdrew. They are listed separately
because their failure mode is different: nobody contradicted them, the
argument was found not to close.

\begin{itemize}
  \item \textbf{The tier collapse as a kinematic degree bound.}
    Proposed with a falsifiable prediction --- that the next shape tier
    would first appear at sixth order --- and retracted eleven days
    later. The count bounded the vertices of $H_4$; the carrier
    projection $d^{\dagger}(\cdot)d$ contributes one further power of
    the Bloch variable that the argument did not carry. Recorded as
    \srcpath{docs/decisions/0005-retracting-the-degree-bound.md}. The
    reformulation stands, the mechanism does not, and the collapse is
    now established by the range-one argument of \S\ref{sec:fourth}
    instead.
  \item \textbf{An unrestricted exponential-source radius.} Inferred from
    independent bounded sources; the inference does not hold, and the
    reformulation is Problem~\ref{prob:g17}.
  \item \textbf{Closure of the source-energy transport budget.} Claimed
    complete, rejected on review: a vacuum-only skew generator does not
    control a source that the generator transports rather than
    annihilates. The corrected partial results are retained and the
    obligation is Problem~\ref{prob:r10}.
\end{itemize}

\noindent The falsified entries above --- the mobility rule, the
stranded-flux backend, the dark centre-only circuits and the global
fixed-window firewall --- are a different category and are not repeated
here: each was killed by a counterexample or an exact computation
rather than by an argument failing to close.

Each of these was, at some point, written down in this program's own
documents as established. That is the reason the machine-verification
tier of \S\ref{sec:howto} is computed rather than assigned, and the
reason the present document quotes its counts from the ledgers by macro
rather than by hand.
